\documentclass[11pt]{article}

\usepackage[a4paper,margin=1in]{geometry}
\usepackage{amsmath,amssymb}
\usepackage{graphicx}
\usepackage{booktabs}
\usepackage{hyperref}
\usepackage{enumitem}

\title{From Keplerian Swarms to Radiatively Supported Bubbles: A Continuum Framework for Dyson Architectures}
\author{Draft in progress}
\date{\today}

\begin{document}

\maketitle

\begin{abstract}
The Dyson Swarm concept is often treated as the most plausible route toward stellar-scale
energy collection, yet its conventional Keplerian formulation contains a geometric
difficulty: large numbers of collectors placed at comparable radii inevitably occupy orbital
planes that intersect at nodal points. Building on established displaced non-Keplerian orbit
(DNKO) theory rather than proposing a fundamentally new family of solar-sail dynamics, this
paper argues that Dyson architectures are more usefully understood as a continuous
support-and-stratification spectrum bridging the Keplerian swarm limit and the fully
radiatively supported bubble/statite limit. Within that spectrum, the low-latitude
heliocentric displaced-orbit regime is reinterpreted as an analytic architecture framework
for layered Dyson-swarm design. For an ideal specular sail, the framework yields the support
curves $\beta_{\min}(\phi)$ and $\sigma_{\max}(\phi)$, which convert the architecture
question into a direct areal-density criterion. The resulting paper is positioned as a
theory-grounded architecture/framework study with illustrative low-latitude feasibility
slices at $0.1^\circ$, $0.5^\circ$, and $1.0^\circ$, rather than as a full engineering
realization, a claim of new orbit-family discovery, or a claim that the displaced-orbit
bridge itself is identified here for the first time.
\end{abstract}

\section{Introduction}

Dyson structures remain a standard thought experiment for stellar-scale energy harvesting,
dating back to Dyson's original argument that advanced civilizations may reprocess a large
fraction of stellar luminosity into distributed artificial collectors \cite{dyson1960infrared}.
Among the usual shell/swarm/bubble taxonomy, the swarm is often treated as the least
structurally extreme option. However, a conventional Keplerian swarm leaves a major
organizational problem unresolved: non-coplanar orbital planes at comparable radii must
intersect at nodes, creating an increasingly severe topological and collision-management
burden as swarm density rises.

Modern Dyson-swarm literature has usually left this issue at a qualitative level, often by
appealing to broad collector distributions or radial nesting without providing a compact
criterion for topological separation at fixed stellar distance. A recent exception is
McInnes (2026), which explicitly notes that collisions in orbiting Dyson swarms could in
principle be reduced using displaced non-Keplerian orbits whose planes can be stacked in
parallel \cite{wright2015gies,mcinnes2025ringworlds,mcinnes2026bubbles}.

This is the Dyson-scale problem that motivates the present paper. The aim is not merely to
identify another solar-sail operating point, but to ask whether the usual Dyson taxonomy is
too coarse. If collector populations can transition continuously from purely orbital support
to partially radiative support and eventually to fully radiatively supported configurations,
then shell/swarm/bubble categories are better understood as regions within a larger support
space rather than as wholly disconnected end states.

That is the specifically Dyson-theoretic claim of the paper. The relevant object is no longer
a single favored megastructure label, but a traversable design space in which support,
stratification, payload margin, and deployment logic vary continuously.

At the opposite extreme, radiatively supported concepts such as statites and Dyson bubbles
replace ordinary orbital support with continuous radiation-pressure support
\cite{forward1991statite,mcinnes1999solarsailing}. The broader solar-sail literature already
contains the relevant astrodynamical ingredients: halo orbits, displaced two-body orbits,
displaced non-Keplerian orbit stability and control, and later synchronous heliocentric
variants \cite{mcinnes1992haloorbits1,mcinnes1992haloorbits2,mcinnes1997displaced,mcinnes1998dnko,quarta2020earthsync,bassetto2024marssync}.

The present paper therefore takes a deliberately conservative novelty posture. Its main
novelty is not a fundamentally new orbit family, nor the first identification of the
displaced-orbit bridge from swarm-collision relief to parallel plane stacking, but the
development of that bridge into an analytic architecture criterion for layered Dyson-swarm
design. More broadly, it proposes that Dyson architectures are better understood as a
continuous support spectrum than as a purely discrete shell/swarm/bubble taxonomy. In that
sense, it should be read as an analytic architecture/framework paper rather than as a new
solar-sail dynamics paper.

\section{Theoretical Framework}

For a collector at fixed heliocentric distance $r$ and latitude $\phi$, moving in a circular
path about the stellar rotation axis with angular velocity $\omega$, define the orbital-rate
ratio
\begin{equation}
\nu \equiv \frac{\omega}{\sqrt{\mu/r^3}}.
\end{equation}

Under the ideal-specular-sail assumption, the payload-optimized branch yields the closed-form
support relation
\begin{equation}
\beta_{\min}(\phi)=\frac{3\sqrt{3}}{2}\sin\phi
\end{equation}
\noindent and the equivalent areal-density limit
\begin{equation}
\sigma_{\max}(\phi)=\frac{2\sigma^*}{3\sqrt{3}\sin\phi}
\end{equation}

so that the architecture question becomes a direct screening criterion:
\begin{equation}
\sigma_{\mathrm{sys}} < \sigma_{\max}(\phi).
\end{equation}

Within the broader DNKO literature, the role of this result in the present paper is not to
claim a new equilibrium family, but to make the Dyson-swarm architecture question analyzable
in closed form. The optimal-angle ingredient behind this result is itself standard prior art
in displaced-orbit analysis, and is used here as an architecture-screening tool rather than
presented as a new derivation. In this way, the support curve acts not only as a local orbit
criterion, but as a parameterization of one tractable low-$\beta$ segment within a larger
Dyson support continuum.

That broader continuum view is essential to the paper's intended contribution. In the present
framing, the key coordinates are not merely orbit labels, but support variables: $\beta$,
$\phi$, $\nu$, and $\sigma_{\max}(\phi)$. These quantities allow one to interpret familiar
Dyson categories as regions within a common support space, rather than as entirely separate
conceptual islands.

One boundary clarification is especially important: the familiar angle
$\phi_c = \arctan(1/\sqrt{2}) \approx 35.264^\circ$ marks the endpoint of the present
payload-optimized displaced-orbit branch, where $\nu \to 0$ and $\beta \to 1.5$. It should
not be read as the terminal latitude of the full Dyson support continuum.

\section{Architecture Reframing}

The intended paper claim is framework-level rather than full-engineering closure, and
architectural rather than dynamical in its novelty posture. The core claim is not merely that
a low-$\beta$ displaced-orbit regime exists, but that Dyson architectures can be usefully
organized as a continuous support spectrum ranging from purely orbital to purely
radiatively-supported configurations. McInnes (2026) already points to the same geometric
bridge in compressed form; what is added here is a full architecture language built around
that bridge.

This is also why the Dyson framing is substantive rather than decorative. The present paper
does not claim that shell, swarm, and bubble architectures are engineeringly interchangeable.
It claims that they can be parameterized within a common support logic, while still occupying
very different regions of cost, control, and payload consequence. The specific role of MDDS is
to take a low-$\beta$ region that recent Dyson literature has signposted only briefly and
make it explicit as a usable framework.

For Dyson theory, this replaces a discourse organized around static end states with one
organized around reachable regions of design space. That is the level at which the paper
claims conceptual novelty; the low-latitude examples that follow remain supporting slices
rather than the paper's principal intellectual center.

\section{Low-Latitude Illustrative Analysis}

The main-text examples are intentionally subordinate to the continuum claim above. Their role
is not to dominate the paper, but to demonstrate that the low-$\beta$ segment of that support
spectrum is non-empty and physically meaningful. The representative slices currently focus on:
\begin{itemize}[leftmargin=1.5em]
  \item $\phi = \theta_\oplus \approx 0.00244^\circ$
  \item $\phi = 0.1^\circ$
  \item $\phi = 0.5^\circ$
  \item $\phi = 1.0^\circ$
\end{itemize}

A useful entry-level characteristic angle is the angular radius of Earth as seen from the Sun,
\[
\theta_\oplus \approx 0.00244^\circ,
\]
for which the present framework gives $\beta_{\min}(\theta_\oplus)\approx 1.11\times10^{-4}$ and
$\sigma_{\max}(\theta_\oplus)\approx 13.83\ \mathrm{kg\,m^{-2}}$, indicating that the pure areal-density
threshold has already entered a regime broadly compatible with present lightweight spacecraft systems.
Geometrically this corresponds to an off-plane separation of about one Earth radius,
$z \approx 6{,}371\ \mathrm{km}$; under the Earth-synchronous slice, the associated inward
radius correction is only about $3.0\times10^3\ \mathrm{km}$.

At the same time, this entry-level observation should not be overextended: the window tightens
rapidly with latitude, and cases such as $0.5^\circ$ and $1.0^\circ$ are better interpreted
as screening-level examples than as near-term engineering claims \cite{macdonald2011applications,mansell2023lightsail2}.

This reading is reinforced by comparison with flown and near-flight sailcraft. LightSail 2
corresponds to a mission-level loading of roughly $156\ \mathrm{g\,m^{-2}}$, while NEA Scout is
of similar order, near $160\ \mathrm{g\,m^{-2}}$ \cite{mansell2023lightsail2,johnson2017neascout}.
These values sit far below the $\theta_\oplus$ threshold and below the $0.1^\circ$ threshold,
but well above the $1^\circ$ threshold, which is exactly why the present paper treats the
former as entry regimes and the latter as a more demanding extrapolation.

\section{Boundaries and Next Steps}

The synchronization-constrained branch should be read as a useful recasting of known
synchronous DNKO design space \cite{heiligers2015earthfollowing,quarta2020earthsync,bassetto2024marssync},
not as an independent orbit-family discovery. Likewise, the paper motivates a reduced
orbital-topology burden for layered swarms, but does not yet claim a fully quantified
collision or conjunction metric.

The topology benefit can nevertheless be stated in a precise geometric sense. A same-radius
Keplerian swarm generates repeated nodal crossings between non-coplanar planes, whereas the
MDDS construction replaces those crossings with layered latitude bands that are explicitly
separated in the normal direction while remaining near the same collection radius. A fuller
future traffic analysis can therefore be posed in terms of node-intersection count, minimum
normal separation, and conjunction-corridor density, but the present manuscript intentionally
stops at the architecture-defining geometry rather than claiming a closed operations metric.

The broader displaced-orbit literature already establishes that stability and station-keeping
are nontrivial questions for such trajectories; existing control results make the architecture
direction plausible, but the present paper does not yet close that layer of analysis
\cite{mcinnes1998dnko,bookless2006thesis,dachwald2005degradation}. The same is true for
optical realism, heavier system-level mass closure, deployment mechanics, and technosignature
modeling. Those questions are natural next steps, but they belong on top of the framework
established here rather than inside its first statement.

\section{Figures}

Current generated figure assets:
\begin{itemize}[leftmargin=1.5em]
  \item \texttt{Paper/figures/concept/fig1\_keplerian\_deadlock.svg}
  \item \texttt{Paper/figures/concept/fig2\_mdds\_stratified\_rings.svg}
  \item \texttt{Paper/figures/results/fig3\_support\_curves.svg}
  \item \texttt{Paper/figures/results/fig4\_low\_latitude\_window.svg}
  \item \texttt{Paper/figures/results/fig5\_sync\_radius\_shift.svg}
\end{itemize}

\bibliographystyle{unsrt}
\bibliography{references/bibliography}

\end{document}
